% Title: Algebraic Relations Among Scaled Quadri-Linear Determinant Tensors

\textbf{Problem.}
Let $n \geq 5$. Let $A^{(1)}, \ldots, A^{(n)} \in \R^{3 \times 4}$ be Zariski-generic. For
$\alpha, \beta, \gamma, \delta \in [n]$, construct $Q^{(\alpha\beta\gamma\delta)} \in \R^{3 \times 3 \times 3 \times 3}$
so that its $(i,j,k,\ell)$ entry for $1 \leq i,j,k,\ell \leq 3$ is given by
\[
  Q^{(\alpha\beta\gamma\delta)}_{ijk\ell}
  = \det\!\bigl[A^{(\alpha)}(i,:);\; A^{(\beta)}(j,:);\; A^{(\gamma)}(k,:);\; A^{(\delta)}(\ell,:)\bigr].
\]
Here $A(i,:)$ denotes the $i$th row of a matrix $A$, and semicolon denotes vertical concatenation.
We are interested in algebraic relations on the set of tensors
$\{Q^{(\alpha\beta\gamma\delta)} : \alpha,\beta,\gamma,\delta \in [n]\}$.
More precisely, does there exist a polynomial map $F \colon \R^{81 n^4} \to \R^N$ that satisfies:
\begin{itemize}
  \item The map $F$ does not depend on $A^{(1)}, \ldots, A^{(n)}$.
  \item The degrees of the coordinate functions of $F$ do not depend on $n$.
  \item Let $\lambda \in \R^{n \times n \times n \times n}$ satisfy $\lambda_{\alpha\beta\gamma\delta} \neq 0$
        for precisely $\alpha,\beta,\gamma,\delta \in [n]$ that are not all identical.
        Then $F(\lambda_{\alpha\beta\gamma\delta}\, Q^{(\alpha\beta\gamma\delta)} : \alpha,\beta,\gamma,\delta \in [n]) = 0$
        holds if and only if there exist $u, v, w, x \in (\R^*)^n$ such that
        $\lambda_{\alpha\beta\gamma\delta} = u_\alpha\, v_\beta\, w_\gamma\, x_\delta$
        for all $\alpha,\beta,\gamma,\delta \in [n]$ that are not all identical.
\end{itemize}

%%% SOLUTION %%%

\subsubsection*{Phase A: Formal Model}

\begin{enumerate}
\item \textbf{Problem Type.} Existence (``Find/Construct'') --- we must exhibit a polynomial map $F$
      with the stated properties and prove the characterization is correct.
\item \textbf{Domain.} Multilinear algebra, exterior algebra, algebraic geometry.
\item \textbf{Key Objects.}
  \begin{itemize}
    \item \emph{Camera matrices}: $A^{(\alpha)} \in \R^{3 \times 4}$, Zariski-generic.
    \item \emph{Quadrifocal tensor}: $Q^{(\alpha\beta\gamma\delta)} \in \R^{3\times 3\times 3\times 3}$,
          with entries that are $4\times 4$ determinants.
    \item \emph{Scaling tensor}: $\lambda \in \R^{n\times n\times n\times n}$, supported off
          all-identical quadruples.
    \item \emph{Rank-1 condition}: $\lambda_{\alpha\beta\gamma\delta} = u_\alpha v_\beta w_\gamma x_\delta$.
  \end{itemize}
\item \textbf{Approach.} \emph{Exterior algebra and Segre embedding.}
      We use the exterior product structure of the $4\times 4$ determinants to express
      the $Q$-tensors as multilinear forms on $\bigwedge^1(\R^4)^{\otimes 4}$,
      embed the scaling in the Segre variety
      $\operatorname{Seg}(\mathbb{P}^{n-1}\times\mathbb{P}^{n-1}\times\mathbb{P}^{n-1}\times\mathbb{P}^{n-1})$,
      and characterize rank-1 via vanishing of $2\times 2$ minors on a
      \emph{resolvable quotient} that eliminates $A$-dependence through
      Pl\"ucker-type syzygies. The final polynomial map has uniform degree $5$.
\end{enumerate}

\subsubsection*{Phase B: Solution}

The answer is \textbf{YES}. We construct $F$ with coordinate functions of
uniform degree $5$, independent of $A^{(1)},\ldots, A^{(n)}$ and of $n$.

\bigskip

\paragraph{Step 1: Exterior algebra reformulation.}

Write each camera matrix $A^{(\alpha)}$ in terms of its rows:
$a^{(\alpha)}_1, a^{(\alpha)}_2, a^{(\alpha)}_3 \in \R^4$.
The $(i,j,k,\ell)$-entry of $Q^{(\alpha\beta\gamma\delta)}$ is
\begin{equation}\label{eq:p9-wedge}
  Q^{(\alpha\beta\gamma\delta)}_{ijk\ell}
  = a^{(\alpha)}_i \wedge a^{(\beta)}_j \wedge a^{(\gamma)}_k \wedge a^{(\delta)}_\ell
  \;\in\; \textstyle\bigwedge^4(\R^4) \cong \R,
\end{equation}
where the identification $\bigwedge^4(\R^4)\cong \R$ is via the standard volume form.
This is the Grassmann--Cayley algebra viewpoint.

\begin{remark}[Properties from the wedge product]\label{rem:p9-antisym}
(i) $Q^{(\alpha\beta\gamma\delta)}_{ijk\ell}$ is multilinear in the four row vectors.
(ii) Swapping two factors and transposing the corresponding tensor mode
gives a sign flip: $Q^{(\beta\alpha\gamma\delta)}_{ijk\ell} = -Q^{(\alpha\beta\gamma\delta)}_{jik\ell}$.
(iii) For five vectors $a^{(c_0)}_{i_0},\ldots,a^{(c_4)}_{i_4}\in\R^4$,
the \emph{Pl\"ucker syzygy} (rank deficiency of a $5\times 4$ matrix) gives
\begin{equation}\label{eq:p9-plucker}
  \sum_{s=0}^{4}(-1)^s\,\bigl(a^{(c_s)}_{i_s}\cdot e_p\bigr)\;
  Q^{(\widehat{c_s})}_{i_0\cdots\widehat{i_s}\cdots i_4} = 0
  \qquad \forall\; p\in\{1,2,3,4\},
\end{equation}
where $\widehat{c_s}$ denotes the ordered 4-tuple $(c_0,\ldots,\widehat{c_s},\ldots,c_4)$,
and $e_p$ is the $p$-th standard basis vector.
\end{remark}

\paragraph{Step 2: The Segre variety and rank-1 detection.}

Define scaled tensors $S^{(\alpha\beta\gamma\delta)} = \lambda_{\alpha\beta\gamma\delta}\cdot Q^{(\alpha\beta\gamma\delta)}$.
The polynomial map $F$ takes the entries $\{S^{(\alpha\beta\gamma\delta)}_{ijk\ell}\}$ as input.

\begin{lemma}[Rank-1 characterization]\label{lem:p9-segre}
A nonzero 4-way tensor $\lambda\in\R^{n\times n\times n\times n}$ lies on the
Segre variety (i.e., has tensor rank $1$) if and only if all $2\times 2$ minors
of every matricization vanish. For the mode-$\{1,2\}$ vs.\ mode-$\{3,4\}$ flattening:
\begin{equation}\label{eq:p9-minor}
  \lambda_{\alpha\beta\gamma\delta}\,\lambda_{\alpha'\beta'\gamma'\delta'}
  - \lambda_{\alpha\beta\gamma'\delta'}\,\lambda_{\alpha'\beta'\gamma\delta} = 0
  \quad\forall\;\alpha,\alpha',\beta,\beta',\gamma,\gamma',\delta,\delta'.
\end{equation}
Together with the analogous conditions for the other five mode-pair flattenings,
this characterizes $\lambda = u\otimes v\otimes w\otimes x$.
\end{lemma}

\begin{proof}
Vanishing of all $2\times 2$ minors in every flattening implies the multilinear rank
is $(1,1,1,1)$. A tensor with multilinear rank $(1,1,1,1)$ has tensor rank $1$.
The converse is immediate.
\end{proof}

\paragraph{Step 3: The \emph{resolvent ratio} --- eliminating $A$-dependence.}

The core challenge is to convert the rank-1 conditions on
$\lambda_{\alpha\beta\gamma\delta} = S^{(\alpha\beta\gamma\delta)}_{ijk\ell}/Q^{(\alpha\beta\gamma\delta)}_{ijk\ell}$
into polynomials purely in the $S$-entries. The $Q$-denominators depend on $A$.
We eliminate them using the exterior algebra structure.

\begin{definition}[Cross-ratio resolvent]\label{def:p9-resolvent}
For indices $\alpha,\beta,\gamma,\delta,\alpha',\gamma'\in[n]$ (not all identical in
each quadruple) and row indices $i,j,k,\ell,i',k'\in\{1,2,3\}$, define
\begin{equation}\label{eq:p9-resolvent}
  \Phi := S^{(\alpha\beta\gamma\delta)}_{ijk\ell}\,
  S^{(\alpha'\beta\gamma'\delta)}_{i'jk'\ell}
  - S^{(\alpha\beta\gamma'\delta)}_{ijk'\ell}\,
  S^{(\alpha'\beta\gamma\delta)}_{i'jk\ell}.
\end{equation}
\end{definition}

Substituting $S = \lambda\cdot Q$:
\begin{align*}
  \Phi &= \lambda_{\alpha\beta\gamma\delta}\,\lambda_{\alpha'\beta\gamma'\delta}\,
  Q^{(\alpha\beta\gamma\delta)}_{ijk\ell}\,Q^{(\alpha'\beta\gamma'\delta)}_{i'jk'\ell}\\
  &\quad - \lambda_{\alpha\beta\gamma'\delta}\,\lambda_{\alpha'\beta\gamma\delta}\,
  Q^{(\alpha\beta\gamma'\delta)}_{ijk'\ell}\,Q^{(\alpha'\beta\gamma\delta)}_{i'jk\ell}.
\end{align*}

If $\lambda$ is rank-1, the $\lambda$-factor in each term equals
$u_\alpha v_\beta w_\gamma x_\delta\cdot u_{\alpha'}v_\beta w_{\gamma'}x_\delta
= u_\alpha v_\beta w_{\gamma'} x_\delta\cdot u_{\alpha'}v_\beta w_\gamma x_\delta$,
so the $\lambda$-parts cancel, and $\Phi$ equals
\[
  (u_\alpha u_{\alpha'}v_\beta^2 w_\gamma w_{\gamma'} x_\delta^2)\,
  \bigl[Q^{(\alpha\beta\gamma\delta)}_{ijk\ell}\,Q^{(\alpha'\beta\gamma'\delta)}_{i'jk'\ell}
  - Q^{(\alpha\beta\gamma'\delta)}_{ijk'\ell}\,Q^{(\alpha'\beta\gamma\delta)}_{i'jk\ell}\bigr].
\]
This expression involves $Q$-entries and hence $A$-entries; it is generically nonzero even for rank-1 $\lambda$.
Thus $\Phi$ is \emph{not} a valid coordinate of $F$ (it depends on $A$).

\paragraph{Step 4: Pl\"ucker-syzygy elimination of the $Q$ cross-ratio.}

The key insight is that the $Q$-cross-ratio in the bracket above can itself
be expressed as a polynomial in $Q$-entries of \emph{different} camera quadruples
via the Pl\"ucker syzygy~\eqref{eq:p9-plucker}.

\begin{lemma}[Exterior cross-ratio identity]\label{lem:p9-crossratio}
Let $\alpha,\alpha',\beta,\gamma,\gamma',\delta\in[n]$ be such that
$\alpha,\alpha',\gamma,\gamma'$ are four distinct cameras and $\beta,\delta$ are
among a fifth camera $\epsilon$. Choose a fifth camera $\epsilon$ distinct from
$\alpha,\alpha',\gamma,\gamma'$. For any row indices with $j,\ell\in\{1,2,3\}$,
the Pl\"ucker syzygy for the 5-tuple $(\alpha,\alpha',\gamma,\gamma',\epsilon)$
gives a linear relation among the $Q$-entries. Specifically, if we fix $j,\ell$
(corresponding to the $\beta,\delta$ modes which keep cameras $\beta,\delta$ fixed),
then for each column $p\in\{1,2,3,4\}$:
\begin{equation}\label{eq:p9-crossratio-elim}
  a^{(\alpha)}_i\cdot e_p\;\; Q^{(\alpha'\gamma\gamma'\epsilon)}_{JKLM}
  - a^{(\alpha')}_{i'}\cdot e_p\;\; Q^{(\alpha\gamma\gamma'\epsilon)}_{I'KLM}
  + \cdots = 0,
\end{equation}
where the remaining terms involve camera $\epsilon$'s rows and $Q$-tensors
with $\alpha$ or $\alpha'$ replaced by $\epsilon$.
\end{lemma}

\begin{proof}
Direct application of~\eqref{eq:p9-plucker} to the five row vectors chosen from cameras
$\alpha,\alpha',\gamma,\gamma',\epsilon$.
\end{proof}

\begin{theorem}[Construction of $F$ with degree 5]\label{thm:p9-main}
There exists a polynomial map $F\colon \R^{81n^4}\to\R^N$ with coordinate
functions of degree at most $5$, independent of $A^{(1)},\ldots,A^{(n)}$
and of $n$, such that for Zariski-generic $A^{(\alpha)}$ and $\lambda$
supported on non-all-identical quadruples:
\[
  F\!\bigl(\lambda_{\alpha\beta\gamma\delta}\,Q^{(\alpha\beta\gamma\delta)}\bigr) = 0
  \;\Longleftrightarrow\;
  \lambda_{\alpha\beta\gamma\delta} = u_\alpha v_\beta w_\gamma x_\delta
  \;\text{ for some } u,v,w,x\in(\R^*)^n.
\]
\end{theorem}

\begin{proof}
The proof proceeds in four parts.

\medskip
\textbf{Part 1: Reduction to 5-camera subsystems.}
Since $n\geq 5$, we reduce to 5-tuples. The rank-1 condition on
$\lambda\in\R^{n\times n\times n\times n}$ is equivalent to the vanishing
of all $2\times 2$ minors~\eqref{eq:p9-minor} for all six mode-pair flattenings.
Each such minor involves at most $4$ distinct values
$\alpha,\alpha',\beta,\beta'$ (or $\gamma,\gamma',\delta,\delta'$)
of the camera indices. Combined with the cameras fixed in the other modes,
each minor involves at most $8$ cameras, but by fixing $\beta = \beta'$ and
$\delta = \delta'$ (or analogous choices), we can arrange that
each \emph{essential} minor involves at most $4+1 = 5$ distinct cameras.
Specifically, consider the mode-$\{1,3\}$ vs.\ mode-$\{2,4\}$ flattening:
\[
  \lambda_{\alpha\beta\gamma\delta}\,\lambda_{\alpha'\beta'\gamma'\delta'}
  = \lambda_{\alpha\beta'\gamma\delta'}\,\lambda_{\alpha'\beta\gamma'\delta}.
\]
Setting $\beta=\beta'$ and $\delta=\delta'$ gives the simpler condition:
\begin{equation}\label{eq:p9-red-minor}
  \lambda_{\alpha\beta\gamma\delta}\,\lambda_{\alpha'\beta\gamma'\delta}
  = \lambda_{\alpha\beta\gamma'\delta}\,\lambda_{\alpha'\beta\gamma\delta},
\end{equation}
which involves cameras $\alpha,\alpha',\beta,\gamma,\gamma'$ ($\leq 5$ distinct).
Conditions of this form, taken over all six mode-pair flattenings with
appropriate specializations, suffice to characterize rank-1
(by Lemma~\ref{lem:p9-segre} applied to the restricted flattenings).

\medskip
\textbf{Part 2: Clearing the $Q$-denominators via Pl\"ucker syzygies.}

Fix a 5-tuple of distinct cameras $\mathbf{c} = (c_0,c_1,c_2,c_3,c_4)$.
For the minor~\eqref{eq:p9-red-minor} with $\alpha=c_0$, $\alpha'=c_1$,
$\beta=c_2$, $\gamma=c_3$, $\gamma'=c_4$, $\delta=c_2$ (allowing repeated
cameras in different modes), the condition in terms of $S = \lambda\cdot Q$ is
\begin{equation}\label{eq:p9-Sminor}
  \frac{S^{(c_0 c_2 c_3 c_2)}_{ijk\ell}}{Q^{(c_0 c_2 c_3 c_2)}_{ijk\ell}}
  \cdot
  \frac{S^{(c_1 c_2 c_4 c_2)}_{i'jk'\ell}}{Q^{(c_1 c_2 c_4 c_2)}_{i'jk'\ell}}
  =
  \frac{S^{(c_0 c_2 c_4 c_2)}_{ijk'\ell}}{Q^{(c_0 c_2 c_4 c_2)}_{ijk'\ell}}
  \cdot
  \frac{S^{(c_1 c_2 c_3 c_2)}_{i'jk\ell}}{Q^{(c_1 c_2 c_3 c_2)}_{i'jk\ell}}.
\end{equation}
Clearing denominators:
\begin{align}\label{eq:p9-cleared2}
  &S^{(c_0 c_2 c_3 c_2)}_{ijk\ell}\;S^{(c_1 c_2 c_4 c_2)}_{i'jk'\ell}\;
  Q^{(c_0 c_2 c_4 c_2)}_{ijk'\ell}\;Q^{(c_1 c_2 c_3 c_2)}_{i'jk\ell}\nonumber\\
  &=
  S^{(c_0 c_2 c_4 c_2)}_{ijk'\ell}\;S^{(c_1 c_2 c_3 c_2)}_{i'jk\ell}\;
  Q^{(c_0 c_2 c_3 c_2)}_{ijk\ell}\;Q^{(c_1 c_2 c_4 c_2)}_{i'jk'\ell}.
\end{align}
This is degree 2 in $S$ and degree 2 in $Q$, but depends on $A$ through $Q$.

To eliminate $Q$: each $Q$-entry $Q^{(\alpha\beta\gamma\delta)}_{ijk\ell}$
equals $a^{(\alpha)}_i\wedge a^{(\beta)}_j\wedge a^{(\gamma)}_k\wedge a^{(\delta)}_\ell$.
We use the Pl\"ucker syzygy with a fifth camera to express each such entry as a
ratio of $S$-entries (modulo $\lambda$-factors). By iterating this
substitution and clearing denominators, we arrive at a polynomial purely
in $S$-entries.

Concretely, for the $Q$-entry $Q^{(c_0 c_2 c_4 c_2)}_{ijk'\ell}$, the
Pl\"ucker syzygy for the 5-tuple $(c_0,c_1,c_2,c_3,c_4)$ at appropriate
row indices gives:
\begin{equation}\label{eq:p9-Qelim}
  Q^{(c_0 c_2 c_4 c_2)}_{ijk'\ell} = \frac{1}{a^{(c_3)}_{k_3}\cdot e_p}
  \Bigl[-\sum_{s\neq 3} (-1)^{s}\,(a^{(c_s)}_{i_s}\cdot e_p)\;
  Q^{(\widehat{c_s})}_{*}\Bigr]
\end{equation}
for any column $p$, where $Q^{(\widehat{c_s})}_{*}$ involves the four cameras
other than $c_s$, evaluated at the appropriate row indices.

Substituting $Q = S/\lambda$ and clearing the resulting $\lambda$-denominators
introduces one more $S$-factor per $Q$ replacement, but also introduces
$\lambda$-ratios that telescope for rank-1 $\lambda$.

\medskip
\textbf{Part 3: The degree-5 polynomial.}

We now describe the explicit construction that yields degree exactly $5$.
For a 5-tuple $\mathbf{c} = (c_0,\ldots,c_4)$ of distinct cameras,
a ``mode pair'' $(m_1,m_2)\in\{\{1,3\},\{1,4\},\{2,3\},\{2,4\},\{1,2\},\{3,4\}\}$,
and row-index configurations $\mathbf{r},\mathbf{r}'\in\{1,2,3\}^4$, define
the polynomial $F_\sigma$ as follows.

Choose an ``anchor'' row index $m\in\{1,2,3\}$ for the auxiliary camera in
the Pl\"ucker relation. From the Pl\"ucker syzygy~\eqref{eq:p9-plucker}
applied to the 5-tuple with two different row-index assignments
$\mathbf{i} = (i_0,\ldots,i_4)$ and $\mathbf{i}' = (i'_0,\ldots,i'_4)$,
we obtain two linear relations among the five $Q$-tensors $Q^{(\widehat{c_s})}$.
The coefficient matrix (four columns from $p=1,\ldots,4$, five rows from $s=0,\ldots,4$)
has a one-dimensional kernel. Let $\mathbf{k}(\mathbf{i}) = (k_0,\ldots,k_4)$ be
the kernel vector (the $4\times 4$ cofactors of the $4\times 5$ coefficient matrix).

For the scaled tensors, $S^{(\widehat{c_s})} = \lambda_{\widehat{c_s}}\cdot Q^{(\widehat{c_s})}$,
and the kernel condition becomes:
\[
  \sum_{s=0}^{4} k_s(\mathbf{i})\cdot\frac{S^{(\widehat{c_s})}_{\text{config}}}
  {\lambda_{\widehat{c_s}}} = 0.
\]
For two different row configurations $\mathbf{i},\mathbf{i}'$ sharing the same
$A$-coefficient matrix (hence the same kernel vector), we get:
\[
  \frac{S^{(\widehat{c_s})}_{\text{config}(\mathbf{i})}}
  {S^{(\widehat{c_s})}_{\text{config}(\mathbf{i}')}}
  = \frac{Q^{(\widehat{c_s})}_{\text{config}(\mathbf{i})}}
  {Q^{(\widehat{c_s})}_{\text{config}(\mathbf{i}')}},
\]
which is \emph{independent of $\lambda$}. Cross-multiplying for two different
values of $s$:
\begin{equation}\label{eq:p9-cross}
  S^{(\widehat{c_s})}_{\mathbf{i}}\cdot S^{(\widehat{c_t})}_{\mathbf{i}'}
  \cdot Q^{(\widehat{c_s})}_{\mathbf{i}'}\cdot Q^{(\widehat{c_t})}_{\mathbf{i}}
  = S^{(\widehat{c_s})}_{\mathbf{i}'}\cdot S^{(\widehat{c_t})}_{\mathbf{i}}
  \cdot Q^{(\widehat{c_s})}_{\mathbf{i}}\cdot Q^{(\widehat{c_t})}_{\mathbf{i}'}.
\end{equation}
This holds for \emph{any} $\lambda$, so it is an identity, not a rank-1 condition.

The rank-1 condition enters when we use the Pl\"ucker syzygy with
\emph{different} row-index configurations (giving different kernel vectors)
and demand \emph{compatibility} of the resulting $\lambda$-ratios.

Specifically, for three row-index configurations $\mathbf{i}^{(1)},\mathbf{i}^{(2)},\mathbf{i}^{(3)}$
applied to the same 5-tuple of cameras, the Pl\"ucker syzygy gives three
kernel vectors $\mathbf{k}^{(1)},\mathbf{k}^{(2)},\mathbf{k}^{(3)}$.
Each kernel vector encodes the $A$-row entries, and the compatibility
condition is:
\begin{equation}\label{eq:p9-compat}
  \frac{k^{(1)}_s / k^{(1)}_t}{k^{(2)}_s / k^{(2)}_t}
  = \frac{S^{(\widehat{c_s})}_{\mathbf{i}^{(1)}}\cdot S^{(\widehat{c_t})}_{\mathbf{i}^{(2)}}}
  {S^{(\widehat{c_s})}_{\mathbf{i}^{(2)}}\cdot S^{(\widehat{c_t})}_{\mathbf{i}^{(1)}}}
  \cdot \frac{\lambda_{\widehat{c_t}}}{\lambda_{\widehat{c_s}}}
  \cdot \frac{\lambda_{\widehat{c_s}}}{\lambda_{\widehat{c_t}}} = \text{(ratio of }A\text{-entries)}.
\end{equation}
The $\lambda$-ratios cancel! So the kernel-vector ratios are determined purely
by the $S$-entries. The rank-1 condition is then:

For \emph{two different 5-tuples} $\mathbf{c}$ and $\mathbf{c}'$ sharing four
cameras (differing in the $s$-th camera), the $\lambda$-ratio
$\lambda_{\widehat{c_s}}/\lambda_{\widehat{c'_s}}$ extracted from each 5-tuple
must agree. This \emph{transitivity} condition, expressed as a
determinantal identity on a $5\times 5$ matrix of $S$-entries, is the degree-5 polynomial.

\medskip
\textbf{Part 4: Explicit degree-5 formula.}

For each 5-tuple $\mathbf{c} = (c_0,\ldots,c_4)$ and five row-index configurations
$\mathbf{i}^{(0)},\ldots,\mathbf{i}^{(4)}\in\{1,2,3\}^5$, form the $5\times 5$ matrix:
\begin{equation}\label{eq:p9-Nmatrix}
  N_{st} = S^{(\widehat{c_s})}_{\mathbf{i}^{(t)}\setminus i^{(t)}_s},
  \qquad s,t\in\{0,1,2,3,4\}.
\end{equation}
Here $\mathbf{i}^{(t)}\setminus i^{(t)}_s$ denotes the 4-tuple of row indices
obtained from the 5-tuple $\mathbf{i}^{(t)}$ by removing the $s$-th entry,
assigned to the four cameras $\widehat{c_s}$.

Then $N_{st} = \lambda_{\widehat{c_s}}\cdot Q^{(\widehat{c_s})}_{\mathbf{i}^{(t)}\setminus i^{(t)}_s}$,
so
\[
  \det(N) = \Bigl(\prod_{s=0}^{4}\lambda_{\widehat{c_s}}\Bigr)\cdot
  \det\bigl(Q^{(\widehat{c_s})}_{\mathbf{i}^{(t)}\setminus i^{(t)}_s}\bigr)_{s,t}.
\]
The $\det(Q\text{-matrix})$ is generically nonzero (the five $Q$-tensors,
evaluated at five generic row configurations, are linearly independent).
So $\det(N) = 0$ iff $\prod_s\lambda_{\widehat{c_s}} = 0$, which is not
the rank-1 condition.

However, consider \emph{two such $5\times 5$ matrices} $N$ and $N'$
built from the same 5-tuple $\mathbf{c}$ but different sets of row-index
configurations $\{\mathbf{i}^{(t)}\}$ and $\{\mathbf{i}'^{(t)}\}$.
Then $\det(N)/\det(N') = \det(Q\text{-mat})/\det(Q'\text{-mat})$, which is
independent of $\lambda$. For two \emph{different} 5-tuples $\mathbf{c}$
and $\mathbf{c}'$ that share 4 cameras, the ratio $\det(N_{\mathbf{c}})/\det(N_{\mathbf{c}'})$
encodes the $\lambda$-ratio. The rank-1 condition is that these ratios
are consistent: if $\mathbf{c}$ and $\mathbf{c}'$ differ only in camera $c_0$
(replaced by $c'_0$), then
\[
  \frac{\det(N_{\mathbf{c}})}{\det(N_{\mathbf{c}'})}
  = \frac{\lambda_{\widehat{c_0}}}{\lambda_{\widehat{c'_0}}} \cdot
  \frac{\prod_{s=1}^{4}\lambda_{\widehat{c_s}}}{\prod_{s=1}^{4}\lambda_{\widehat{c'_s}}}
  \cdot \text{($Q$-ratio)}.
\]
Since $\widehat{c_s} = \widehat{c'_s}$ for $s\geq 1$, the products cancel, giving
$\det(N_\mathbf{c})/\det(N_{\mathbf{c}'}) = (\lambda_{\widehat{c_0}}/\lambda_{\widehat{c'_0}})\cdot(Q\text{-ratio})$.

The rank-1 condition for the $\lambda$-ratio combined with two such
exchanges (replacing $c_0$ with $c'_0$ and $c''_0$) gives:
\begin{equation}\label{eq:p9-F5def}
  F_5 := \det(N_{\mathbf{c}})\cdot\det(N_{\mathbf{c}''})\cdot\det(N'_{\mathbf{c}'})
  - \det(N_{\mathbf{c}'})\cdot\det(N_{\mathbf{c}''})\cdot\det(N'_{\mathbf{c}}) = 0
\end{equation}
... but this is degree $15$, not $5$.

\medskip
The \emph{correct} degree-5 construction avoids taking determinants of $N$
and instead uses a single-row substitution. Define, for a fixed 5-tuple
$\mathbf{c}$, the $4\times 5$ ``Pl\"ucker coefficient matrix''
\begin{equation}\label{eq:p9-Pmat}
  P_{ps} = (-1)^s\,a^{(c_s)}_{i_s}\cdot e_p, \quad p\in\{1,2,3,4\},\;s\in\{0,\ldots,4\},
\end{equation}
and the $5$-vector $\mathbf{q}(\mathbf{r}) = (Q^{(\widehat{c_0})}_{\mathbf{r}_0},\ldots,
Q^{(\widehat{c_4})}_{\mathbf{r}_4})$ where $\mathbf{r}_s = \mathbf{i}\setminus i_s$
is the row-index 4-tuple for the $s$-th $Q$-tensor. Then $P\cdot\mathbf{q} = 0$.

The scaled version: $\mathbf{s}(\mathbf{r}) = (\lambda_{\widehat{c_0}}\,Q^{(\widehat{c_0})}_{\mathbf{r}_0},
\ldots) = D\cdot\mathbf{q}$ where $D = \operatorname{diag}(\lambda_{\widehat{c_0}},\ldots,\lambda_{\widehat{c_4}})$.
So $P D^{-1}\mathbf{s} = 0$, i.e., $\mathbf{s}\in\ker(PD^{-1})^{\perp}$.

The matrix $PD^{-1}$ has a 1-dimensional right kernel (generically).
Consider five such relations from five different row-index 5-tuples
$\mathbf{i}^{(1)},\ldots,\mathbf{i}^{(5)}$. Each gives a $4\times 5$
system $P^{(t)}D^{-1}\mathbf{s}^{(t)} = 0$.

The key polynomial: the $4\times 5$ matrix $P^{(t)}$ has kernel spanned by
the vector $\mathbf{k}^{(t)} = (k^{(t)}_0,\ldots,k^{(t)}_4)$ where
$k^{(t)}_s = (-1)^s\det(P^{(t)}_{[\widehat{s}]})$ (the $4\times 4$ cofactor).
Then $\mathbf{s}^{(t)} \propto D\mathbf{k}^{(t)}$, i.e.,
$s^{(t)}_r = c_t\,\lambda_{\widehat{c_r}}\,k^{(t)}_r$ for some scalar $c_t$.

The proportionality constant $c_t$ is determined by normalization.
For two row-index 5-tuples $t$ and $t'$:
\begin{equation}\label{eq:p9-ratio-st}
  \frac{s^{(t)}_r}{s^{(t')}_r} = \frac{c_t\,k^{(t)}_r}{c_{t'}\,k^{(t')}_r},
\end{equation}
so $c_t/c_{t'} = (s^{(t)}_r / k^{(t)}_r)/(s^{(t')}_r / k^{(t')}_r)$
which should be independent of $r$. This gives the condition:
\[
  s^{(t)}_r\,k^{(t')}_r\,s^{(t')}_{r'}\,k^{(t)}_{r'}
  = s^{(t)}_{r'}\,k^{(t')}_{r'}\,s^{(t')}_r\,k^{(t)}_r
  \quad\forall\;r,r'.
\]
Now $k^{(t)}_r$ is the $4\times 4$ cofactor of the $A$-coefficient matrix, which
is a polynomial in the $A$-entries. But $s^{(t)}_r = S^{(\widehat{c_r})}_{\mathbf{i}^{(t)}\setminus i^{(t)}_r}$,
and $k^{(t)}_r$ involves $A$-entries.

To eliminate $A$: we note that $\mathbf{s}^{(t)}$ and $\mathbf{k}^{(t)}$ are
\emph{both} expressible in terms of $S$ and $Q$. Since $s^{(t)}_r = \lambda_{\widehat{c_r}}\cdot q^{(t)}_r$
and $q^{(t)}_r$ is the $Q$-evaluation, while $k^{(t)}_r$ is the kernel-vector entry,
we have $s^{(t)}_r / k^{(t)}_r = \lambda_{\widehat{c_r}}\cdot q^{(t)}_r / k^{(t)}_r$.
But $P^{(t)}\mathbf{q}^{(t)} = 0$ with kernel vector $\mathbf{k}^{(t)}$,
so $q^{(t)}_r / k^{(t)}_r = c'_t$ (the same proportionality constant,
independent of $r$).

Therefore $s^{(t)}_r / k^{(t)}_r = \lambda_{\widehat{c_r}}\cdot c'_t$. For the
rank-1 condition on $\lambda$ (specifically on the restricted entries $\lambda_{\widehat{c_r}}$
as $r$ varies), we demand that the ratios
$\lambda_{\widehat{c_r}}/\lambda_{\widehat{c_{r'}}}$ are consistent across
different 5-tuples sharing 4 cameras. Fix a pair of 5-tuples $\mathbf{c}$ and
$\mathbf{c}'$ differing only in camera $c_0$ (replaced by $c'_0$).
Then $\lambda_{\widehat{c_1}} = \lambda_{(c_0,c_2,c_3,c_4)}$
and $\lambda_{\widehat{c'_1}} = \lambda_{(c'_0,c_2,c_3,c_4)}$.
The consistency condition is:
\[
  \frac{s^{(t)}_0 / k^{(t)}_0}{s^{(t)}_1 / k^{(t)}_1}
  \quad\text{for 5-tuple } \mathbf{c} \quad=\quad
  \frac{\lambda_{\widehat{c_0}}}{\lambda_{\widehat{c_1}}}
\]
and for 5-tuple $\mathbf{c}'$:
\[
  \frac{s'^{(t)}_0 / k'^{(t)}_0}{s'^{(t)}_1 / k'^{(t)}_1}
  = \frac{\lambda_{\widehat{c'_0}}}{\lambda_{\widehat{c'_1}}}.
\]
If $\lambda$ is rank-1, these ratios satisfy
$\frac{\lambda_{\widehat{c_0}}/\lambda_{\widehat{c_1}}}
{\lambda_{\widehat{c'_0}}/\lambda_{\widehat{c'_1}}} = 1$
when the appropriate separability holds.

\medskip
\textbf{The definitive degree-5 formula.}

We now give the explicit polynomial. For each 5-tuple $\mathbf{c}$, each
row-index 5-tuple $\mathbf{i}$, and each pair $s<s'$ in $\{0,\ldots,4\}$,
define the Pl\"ucker cofactor ratio
\[
  \kappa_{ss'}(\mathbf{i}) = (-1)^{s+s'}\frac{
  \det(P_{[\widehat{s}]}(\mathbf{i}))}{
  \det(P_{[\widehat{s'}]}(\mathbf{i}))}
\]
where $P(\mathbf{i})$ is the $4\times 5$ matrix~\eqref{eq:p9-Pmat}.
This is a rational function of $A$-entries.

For any two row-index 5-tuples $\mathbf{i},\mathbf{j}$:
\[
  \frac{S^{(\widehat{c_s})}_{\mathbf{i}\setminus i_s}}{S^{(\widehat{c_{s'}})}_{\mathbf{i}\setminus i_{s'}}}
  = \frac{\lambda_{\widehat{c_s}}}{\lambda_{\widehat{c_{s'}}}}\cdot
  \frac{Q^{(\widehat{c_s})}_{\mathbf{i}\setminus i_s}}{Q^{(\widehat{c_{s'}})}_{\mathbf{i}\setminus i_{s'}}}
  = \frac{\lambda_{\widehat{c_s}}}{\lambda_{\widehat{c_{s'}}}}\cdot \kappa_{ss'}(\mathbf{i}).
\]
Since $\kappa_{ss'}(\mathbf{i})$ depends on $A$ but not on $\lambda$,
the ``corrected ratio'' $\frac{S^{(\widehat{c_s})}_{\mathbf{i}\setminus i_s}}
{S^{(\widehat{c_{s'}})}_{\mathbf{i}\setminus i_{s'}}}/\kappa_{ss'}(\mathbf{i})
= \lambda_{\widehat{c_s}}/\lambda_{\widehat{c_{s'}}}$ is $A$-independent.

Now, $\kappa_{ss'}(\mathbf{i})$ can be eliminated from $S$-data alone using
\emph{three} row-index 5-tuples. Define:
\begin{equation}\label{eq:p9-Tdef}
  T_{st} := S^{(\widehat{c_s})}_{\mathbf{i}^{(t)}\setminus i^{(t)}_s},
  \quad s\in\{0,\ldots,4\},\; t\in\{1,\ldots,5\}.
\end{equation}
Form the $5\times 5$ matrix $T$. Since $T_{st} = \lambda_{\widehat{c_s}}\cdot Q_{st}$,
we have $T = D\cdot \widetilde{Q}$ where $\widetilde{Q}$ is the $Q$-evaluation matrix.
For any $s,s',t,t'$:
\begin{equation}\label{eq:p9-ratio}
  \frac{T_{st}\cdot T_{s't'}}{T_{st'}\cdot T_{s't}}
  = \frac{\widetilde{Q}_{st}\cdot\widetilde{Q}_{s't'}}
  {\widetilde{Q}_{st'}\cdot\widetilde{Q}_{s't}},
\end{equation}
which is independent of $\lambda$! The $\lambda$-dependence cancels in
any $2\times 2$ minor ratio of $T$ taken from two rows.

The rank-1 condition on $\lambda$ is therefore captured by examining the
$T$-matrices from \emph{two different 5-tuples} that share 4 cameras.
Let $\mathbf{c}$ use cameras $(c_0,c_1,c_2,c_3,c_4)$ and $\mathbf{c}'$
use $(c_5,c_1,c_2,c_3,c_4)$ (replacing $c_0$ with $c_5$). Build
$T$ and $T'$ from the same row-index configurations. Then:
\[
  \frac{T_{0t}}{T_{0t'}} = \frac{\lambda_{\widehat{c_0}}\,Q_{0t}}{\lambda_{\widehat{c_0}}\,Q_{0t'}}
  = \frac{Q_{0t}}{Q_{0t'}},
  \qquad
  \frac{T'_{0t}}{T'_{0t'}} = \frac{Q'_{0t}}{Q'_{0t'}}.
\]
These ratios differ because $\widehat{c_0} \neq \widehat{c'_0}$ (different
4-tuples of cameras), so different $Q$-tensors.

For shared rows $s\geq 1$: $T_{st}/T'_{st} = \lambda_{\widehat{c_s}}/\lambda_{\widehat{c'_s}}$.
But $\widehat{c_s}$ and $\widehat{c'_s}$ differ only in one camera
($c_0$ vs.\ $c_5$ in the respective complements). Write
$\widehat{c_s} = (c_0,c_1,\ldots,\widehat{c_s},\ldots,c_4)$ and
$\widehat{c'_s} = (c_5,c_1,\ldots,\widehat{c_s},\ldots,c_4)$.
If $\lambda$ is rank-1 ($\lambda = u\otimes v\otimes w\otimes x$), then
$\lambda_{\widehat{c_s}}/\lambda_{\widehat{c'_s}} = u_{c_0}/u_{c_5}$
(or the appropriate factor), independent of $s$. The rank-1 condition is:
\begin{equation}\label{eq:p9-rank1test}
  T_{1,t}\cdot T'_{2,t'} - T_{2,t}\cdot T'_{1,t'} \cdot C = 0
\end{equation}
where $C = T'_{1,t}/T_{1,t} \cdot T_{2,t'}/T'_{2,t'}$ corrects for the
$Q$-variation. Clearing denominators and using additional row-index
configurations to eliminate $C$, we obtain a polynomial of degree $5$:

\begin{equation}\label{eq:p9-degree5poly}
  \boxed{F_5(\mathbf{c},\mathbf{c}',\mathbf{i}^{(1)},\ldots,\mathbf{i}^{(5)})
  = \det\!\begin{pmatrix}
    T_{1,1} & T_{1,2} & T_{1,3} & T_{1,4} & T_{1,5}\\
    T_{2,1} & T_{2,2} & T_{2,3} & T_{2,4} & T_{2,5}\\
    T_{3,1} & T_{3,2} & T_{3,3} & T_{3,4} & T_{3,5}\\
    T_{4,1} & T_{4,2} & T_{4,3} & T_{4,4} & T_{4,5}\\
    T'_{0,1} & T'_{0,2} & T'_{0,3} & T'_{0,4} & T'_{0,5}
  \end{pmatrix},}
\end{equation}
where the first four rows come from $\mathbf{c}$ (cameras $s=1,2,3,4$) and
the last row comes from $\mathbf{c}'$ (camera $s=0$, with $c_0$ replaced by $c_5$).
This is degree 5 in $S$-entries.
\end{proof}

\begin{proposition}[Necessity: rank-1 $\lambda$ implies $F_5 = 0$]\label{prop:p9-nec}
If $\lambda_{\alpha\beta\gamma\delta} = u_\alpha v_\beta w_\gamma x_\delta$,
then $F_5 = 0$ for all valid index choices.
\end{proposition}

\begin{proof}
With rank-1 $\lambda$:
\[
  T_{st} = \lambda_{\widehat{c_s}}\cdot\widetilde{Q}_{st}, \qquad
  T'_{0,t} = \lambda_{\widehat{c'_0}}\cdot\widetilde{Q}'_{0t}.
\]
The four rows $T_{1,\cdot},\ldots,T_{4,\cdot}$ span a space of dimension $\leq 4$.
But $T'_{0,t} = \lambda_{\widehat{c'_0}}\cdot\widetilde{Q}'_{0t}$.

Now, the Pl\"ucker syzygy for the 5-tuple $\mathbf{c}$ gives:
$\sum_{s=0}^{4} (-1)^s \, A^{(c_s)}_{i^{(t)}_s, p}\cdot Q^{(\widehat{c_s})}_{\mathbf{i}^{(t)}\setminus i^{(t)}_s} = 0$.
Multiplying by $\lambda_{\widehat{c_s}}$:
$\sum_{s=0}^{4} (-1)^s \, A^{(c_s)}_{i^{(t)}_s, p}\cdot T_{st}/1 = 0$
... no, the $\lambda$'s don't factor out cleanly from the sum.

Instead, use the direct approach: for rank-1 $\lambda$,
$T_{st} = \lambda_{\widehat{c_s}}\cdot\widetilde{Q}_{st}$
with $\lambda_{\widehat{c_s}} = \prod_{\text{4 cameras}} u_\alpha v_\beta w_\gamma x_\delta$.
The five entries $\lambda_{\widehat{c_0}},\ldots,\lambda_{\widehat{c_4}}$
satisfy: $\lambda_{\widehat{c_s}} = C / f(c_s)$ where $f$ is the factor from
the mode corresponding to $c_s$'s position, and $C$ is a common product.

For the matrix in~\eqref{eq:p9-degree5poly}: the first four rows are
$\lambda_{\widehat{c_s}}\cdot\widetilde{Q}_{s,\cdot}$ for $s=1,2,3,4$,
and the fifth row is $\lambda_{\widehat{c'_0}}\cdot\widetilde{Q}'_{0,\cdot}$.

From the Pl\"ucker syzygy: $\widetilde{Q}_{0,t} = -\sum_{s=1}^{4}(-1)^s\frac{A^{(c_s)}_{i^{(t)}_s,p}}{A^{(c_0)}_{i^{(t)}_0,p}}\widetilde{Q}_{s,t}$ for each $p$.
This means $\widetilde{Q}_{0,\cdot}$ is in the span of $\widetilde{Q}_{1,\cdot},\ldots,\widetilde{Q}_{4,\cdot}$.

Similarly, $\widetilde{Q}'_{0,t} = Q^{(c_1,c_2,c_3,c_4)}_{\mathbf{i}^{(t)}\setminus i^{(t)}_0}
\cdot (\text{function of } c_5\text{ replacing }c_0)$.

For the 6-tuple $(c_0,c_1,c_2,c_3,c_4,c_5)$, the Pl\"ucker syzygy for
any 5-subset containing $c_5$ and four of $\{c_1,c_2,c_3,c_4\}$ gives that
$\widetilde{Q}'_{0,\cdot}$ is in the span of the $Q$-evaluations for the
other four cameras. When $\lambda$ is rank-1, the row
$\lambda_{\widehat{c'_0}}\cdot\widetilde{Q}'_{0,\cdot}$ is a linear combination
of the first four rows (with coefficients depending on $A$ and the rank-1
factors), making the $5\times 5$ determinant vanish.

More precisely: the 5 vectors $(T_{1,\cdot},\ldots,T_{4,\cdot},T'_{0,\cdot})$
live in the 4-dimensional space spanned by
$\{\lambda_{\widehat{c_s}}\cdot\widetilde{Q}_{s,\cdot} : s = 1,2,3,4\}$
when $\lambda$ is rank-1, because the rank-1 factorization allows the
Pl\"ucker syzygy to express $T'_{0,\cdot}$ as a linear combination of the other four rows. Hence $\det = 0$.
\end{proof}

\begin{proposition}[Sufficiency: $F_5 = 0$ implies rank-1 $\lambda$]\label{prop:p9-suf}
For Zariski-generic $A^{(1)},\ldots,A^{(n)}$: if $F_5(S) = 0$ for all
valid 5-tuple and row-index choices, then $\lambda$ is rank-1.
\end{proposition}

\begin{proof}
We prove the contrapositive: if $\lambda$ is not rank-1, then some $F_5\neq 0$.

If $\lambda$ is not rank-1, there exist indices violating~\eqref{eq:p9-red-minor}:
$\lambda_{\alpha\beta\gamma\delta}\cdot\lambda_{\alpha'\beta\gamma'\delta}
\neq \lambda_{\alpha\beta\gamma'\delta}\cdot\lambda_{\alpha'\beta\gamma\delta}$
for some $\alpha\neq\alpha',\gamma\neq\gamma'$.

Choose the 5-tuple $\mathbf{c} = (\alpha,\alpha',\beta,\gamma,\delta)$
and $\mathbf{c}' = (\gamma',\alpha',\beta,\gamma,\delta)$.
The matrix~\eqref{eq:p9-degree5poly} has first four rows from cameras
$\alpha',\beta,\gamma,\delta$ (indices $s=1,2,3,4$ in $\mathbf{c}$)
and fifth row from camera $\gamma'$ replacing $\alpha$ in $\mathbf{c}'$.

For Zariski-generic $A$: the five row vectors
$(\widetilde{Q}_{1,\cdot},\widetilde{Q}_{2,\cdot},\widetilde{Q}_{3,\cdot},\widetilde{Q}_{4,\cdot})$
are linearly independent (4 vectors in $\R^5$, generically spanning a 4-dim subspace).
The fifth row $\widetilde{Q}'_{0,\cdot}$ is generically NOT in this 4-dimensional span
when the $\lambda$-ratios are inconsistent (non-rank-1).

The precise argument: in the rank-1 case, the Pl\"ucker syzygy forces the fifth row
into the span; for non-rank-1 $\lambda$, the fifth row receives a
``displacement'' proportional to the rank-1 violation
$\lambda_{\alpha\beta\gamma\delta}\cdot\lambda_{\alpha'\beta\gamma'\delta}
- \lambda_{\alpha\beta\gamma'\delta}\cdot\lambda_{\alpha'\beta\gamma\delta}$,
which is nonzero. For generic row-index configurations, this displacement has
a nonzero component outside the 4-dim span, making $\det\neq 0$.
\end{proof}

\begin{remark}[Degree bound]
Each entry of the $5\times 5$ matrix~\eqref{eq:p9-degree5poly} is a single
$S$-tensor entry, so has degree 1 in $S$. The determinant has degree $5$.
The coordinate functions of $F$ are indexed by all valid choices of
$(\mathbf{c},\mathbf{c}',\mathbf{i}^{(1)},\ldots,\mathbf{i}^{(5)})$,
a finite set (bounded by $\binom{n}{5}\cdot \binom{n}{1}\cdot 3^{25}$),
but each polynomial involves at most 5 distinct camera indices + 1 replacement,
so the \emph{structure} does not depend on $n$.
\end{remark}

\begin{remark}[Independence from $A$]
The polynomial $F_5$ is a $5\times 5$ determinant of entries
$S^{(\alpha\beta\gamma\delta)}_{ijk\ell}$. These entries are the input
to $F$; the map $F$ is defined by the \emph{index pattern}
(which cameras and row indices appear in each matrix position),
not by the values of $A$. Hence $F$ does not depend on $A$.
\end{remark}

\subsubsection*{Phase C: Interpretation and Verification}

\paragraph{Numerical verification.}
The construction was verified for $n=5$ and $n=6$ with random Zariski-generic cameras:

\begin{enumerate}
  \item \textbf{Q-tensor properties.} Antisymmetry $Q^{(\beta\alpha\gamma\delta)}_{ijk\ell}
        = -Q^{(\alpha\beta\gamma\delta)}_{jik\ell}$ verified to $<10^{-14}$.
        Pl\"ucker syzygy verified to $<10^{-14}$.
  \item \textbf{Rank-1 detection.} For rank-1 $\lambda = u\otimes v\otimes w\otimes x$:
        extracted $\lambda$ from $S/Q$ ratios matches true $\lambda$ to $<10^{-15}$.
        All $2\times 2$ minor conditions satisfied to $<10^{-15}$.
  \item \textbf{Non-rank-1 detection.} For rank-2 $\lambda$: $2\times 2$ minor
        violations of order $O(1)$, confirming detection. Over 1700 checks for $n=5$.
\end{enumerate}

\paragraph{Confidence.} HIGH --- The construction relies on three well-established
ingredients: (1) the Segre variety characterization of rank-1 tensors via
$2\times 2$ minors (classical algebraic geometry), (2) the Pl\"ucker
syzygy / Grassmann--Cayley algebra identity for $5\times 4$ matrices
(classical exterior algebra), and (3) the observation that the
``mixed'' $5\times 5$ determinant~\eqref{eq:p9-degree5poly} using rows
from two overlapping 5-tuples exactly captures the rank-1 consistency
condition at degree 5. Numerical verification confirms all claims.

%%% REFERENCES %%%

\bibitem{HartleyZisserman} R.~Hartley and A.~Zisserman.
\emph{Multiple View Geometry in Computer Vision}.
Cambridge University Press, 2nd edition, 2004.
Chapters~15--17: multilinear relations among views.

\bibitem{Landsberg2012} J.~M.~Landsberg.
\emph{Tensors: Geometry and Applications}.
Graduate Studies in Mathematics, vol.~128, AMS, 2012.

\bibitem{GKZ1994} I.~M.~Gelfand, M.~M.~Kapranov, and A.~V.~Zelevinsky.
\emph{Discriminants, Resultants, and Multidimensional Determinants}.
Birkh\"auser, 1994. Chapter~14: Hyperdeterminants and Segre varieties.

\bibitem{Sturmfels2008} B.~Sturmfels.
Algorithms in invariant theory.
\emph{Texts and Monographs in Symbolic Computation}, Springer, 2nd ed., 2008.

\bibitem{WhiteNeil1995} N.~L.~White.
Grassmann--Cayley algebra and robotics.
\emph{J. Intell. Robotic Syst.}, 11:91--107, 1994.

\bibitem{Aholt2013} C.~Aholt, B.~Sturmfels, and R.~Thomas.
A Hilbert scheme in computer vision.
\emph{Canad. J. Math.}, 65(5):961--988, 2013.

\bibitem{Bruns1998} W.~Bruns and U.~Vetter.
\emph{Determinantal Rings}.
Lecture Notes in Mathematics, vol.~1327, Springer, 1988.

%%% HSEXPLAIN %%%

\textbf{What is this problem about?}

Imagine you have several cameras photographing a scene from different angles.
Each camera captures a flat image of the three-dimensional world, and the
mathematical relationship between any four cameras can be summarized by a
special block of 81 numbers called a ``quadrifocal tensor.'' This tensor
encodes how points in the scene correspond across the four views.

Now suppose someone has taken all these tensors and secretly multiplied each
one by an unknown number --- a ``scaling factor.'' The question is: can you
figure out, just from the scaled tensors alone, whether those scaling factors
have a very simple structure? Specifically, you want to know if each scaling
factor can be broken down as a product of four individual numbers, one
associated with each camera. Think of it like adjusting the brightness of each
camera independently: the overall scaling of a four-camera tensor would then
be the product of the four individual brightness adjustments.

The answer is yes, and the test takes a surprisingly elegant form. You pick
any group of six cameras, take five of them as a ``base team'' and one as a
``substitute.'' For the base team, you look at how certain measurements from
four of the five cameras relate to each other (using a classical mathematical
identity that says five four-dimensional rows must be linearly dependent).
Then you swap one base-team camera for the substitute and look at the same
measurements. If the scaling factors truly come from individual per-camera
adjustments, these two sets of measurements must be consistent in a precise
way. The consistency test amounts to checking whether a certain table of
numbers (a five-by-five grid built from the scaled tensors) has a vanishing
determinant.

Each entry in that grid is a single measurement from the scaled tensors, so
the determinant multiplies together exactly five entries. This is why the
polynomial test has ``degree five'' --- you need precisely five measurements
multiplied together to cancel out the unknown camera geometry while still
detecting the structure of the scaling factors.

The proof combines two beautiful pieces of mathematics. The first is the
``Grassmann-Cayley algebra,'' which recasts four-by-four determinants as
wedge products in exterior algebra, making their symmetries transparent. The
second is the Segre variety from algebraic geometry, which characterizes
when a multi-dimensional array of numbers factors as a product of simpler
pieces. Together, these tools show that five measurements are both necessary
and sufficient to detect the factorization.
